% Single-file LaTeX source — paste directly into Overleaf (article class).
\documentclass[11pt]{article}

\usepackage[margin=1in]{geometry}
\usepackage{amsmath, amssymb, amsthm}
\usepackage{booktabs}
\usepackage{microtype}
\usepackage{xcolor}
\usepackage{listings}
\usepackage{graphicx}
\usepackage{enumitem}
\usepackage[hidelinks]{hyperref}
\usepackage{url}
\usepackage{caption}
\usepackage{pgfplots}
\pgfplotsset{compat=1.18}
\usepackage{array}

% --- TLA+ listing style -----------------------------------------------------
\definecolor{kwblue}{RGB}{0,0,180}
\definecolor{cmgray}{RGB}{120,120,120}
\definecolor{strgreen}{RGB}{0,120,0}
\lstdefinelanguage{TLA}{
  morekeywords={MODULE, EXTENDS, VARIABLES, CONSTANT, CONSTANTS, ASSUME,
                THEOREM, LEMMA, PROOF, BY, DEF, DEFS, EXCEPT, IF, THEN, ELSE,
                LET, IN, CASE, OTHER, INSTANCE, LOCAL, UNCHANGED,
                ENABLED, CHOOSE, SUBSET, UNION, DOMAIN, TRUE, FALSE, BOOLEAN,
                Init, Next, Spec, TypeOK, Inv},
  sensitive=true,
  morecomment=[l]{\\*},
  morecomment=[s]{(*}{*)},
  morestring=[b]",
}
\lstset{
  basicstyle=\ttfamily\footnotesize,
  keywordstyle=\color{kwblue}\bfseries,
  commentstyle=\color{cmgray}\itshape,
  stringstyle=\color{strgreen},
  columns=fullflexible,
  showstringspaces=false,
  breaklines=true,
  frame=single,
  framerule=0.3pt,
  framesep=4pt,
  xleftmargin=6pt,
  xrightmargin=6pt,
  aboveskip=6pt,
  belowskip=6pt,
}

% --- Title ------------------------------------------------------------------
\title{Fine-Tuning a 20B Open Model for TLA\textsuperscript{+}\\
A Year of Failed Optimisations and What Finally Worked}
\author{Eric Spencer\\
\small \textit{(Loyola University Chicago)}\\
\small \texttt{REDACTED-USER@luc.edu}}
\date{April 7, 2026 (revised)}

\begin{document}
\maketitle

% ---------------------------------------------------------------------------
\begin{abstract}
We report on a year of fine-tuning OpenAI's open-weight \texttt{gpt-oss:20b}
model for TLA\textsuperscript{+} specification generation. Across more than
two hundred reinforcement-learning cycles, four supervised fine-tuning
regimes (broad SFT, incremental SFT, KTO, DPO), an automated rebuild,
deploy, and rollback pipeline, $484$ verifier-labelled
training examples in the augmented archive (of which $73$ pass the
Diamond gate), and a follow-on round of $200$ purpose-generated specs
across ten algorithm families that pass the gate at a $200/200$
yield, the trajectory of the held-out 20-problem benchmark is
striking for how much it does \emph{not} reward the things we did. Our
strongest end-to-end checkpoint reaches $9/20$ SANY and $5/20$ TLC; our
weakest catastrophically forgets the base model and falls to $0/20$ on
both. The interesting structure is in the failures: every time we trained
broadly on a small corpus from scratch the loop knowledge collapsed; every
time we used preference optimisation against a binary verifier reward, the
model drifted toward semantically vacuous specifications. The eventual win
came from inverting the usual pipeline---using verifiers at inference time
through decomposition rather than at training time through preference
loss---and reserving fine-tuning for narrow, rejection-sampled, semantically
audited data. This paper documents the architecture of the training stack
(\S\ref{sec:arch}), the chronological record of the fine-tuning runs and
their benchmark deltas (\S\ref{sec:runs}), the four lessons that the
log files seem to teach (\S\ref{sec:lessons}), and the current
Diamond-filtered SFT recipe (\S\ref{sec:diamond-sft}). It is the
companion to our methods paper on the Diamond gate and piecewise
generation~\cite{spencer2026chattla}, which should be read first; here
we give the engineering side of the same story.
\end{abstract}

% ---------------------------------------------------------------------------
\section{Introduction}

TLA\textsuperscript{+}~\cite{lamport2002specifying} sits in an awkward
place for language modelling. The corpus is small (a few hundred public
specifications), the verifier (SANY plus the TLC model checker) is
deterministic and cheap, and the cost of getting a specification subtly
wrong is high enough that there is no obvious gradient between ``almost
right'' and ``right.'' These properties make it tempting to try every
fine-tuning trick in the contemporary toolbox: broad SFT to absorb the
public corpus, preference optimisation (DPO, KTO) against the verifier,
reinforcement learning from verifier feedback, distillation from larger
teachers. Over the course of a year we tried all of those. Most of them
made the model worse.

We ran the entire experiment against a single underlying checkpoint,
OpenAI's open-weight \texttt{gpt-oss:20b}, served locally through Ollama
and fine-tuned with LoRA adapters merged into BF16 weights for evaluation.
Every benchmark number reported here comes from the same 20-problem
held-out suite (\texttt{BM001}--\texttt{BM020}) that drives our methods
paper~\cite{spencer2026chattla}; every retrain, every merge, every GGUF
build, and every Hugging Face publish is automated and recorded under
\texttt{outputs/}. The point of this paper is not to report a single
headline number---we already have one---but to make the audit trail
visible: which fine-tuning runs helped, which hurt, and what the failures
seem to imply about the shape of the problem.

This revision (April 7, 2026) adds \S\ref{sec:diamond-gen-v2}, which
documents a second-generation Diamond corpus produced by topic-conditioned
multi-agent generation rather than by post-hoc rejection sampling against
a teacher. The new method shifts the binding constraint from
\emph{teacher access} to \emph{topic curriculum coverage}, raises the
gate-pass yield from the $\sim\!15\%$ achievable under rejection sampling
to $200/200$ on its first run, and is the substrate for the
\texttt{diamond-gen-v2} fine-tune now in flight against the same
$20$-problem benchmark and a separately-carved $30$-problem held-out
diamond eval.

The summary, before any of the detail: between version~6 and version~14
of the model, the held-out TLC pass count moved from $1/20$ to $5/20$ to
$0/20$ and back, depending on which lever we pulled. The single
intervention with the largest positive effect on TLC was not a training
change at all; it was a five-step inference-time decomposition of the
generation task. The single intervention with the largest negative effect
was a from-scratch SFT pass on $234$ examples that erased the model's
TLA\textsuperscript{+} prior in a single $15$-epoch run. Both observations
are consistent with the working hypothesis that for low-resource formal
languages, the shape of the verifier reward matters more than the
quantity of supervised data.

We make three contributions in this paper:
\begin{enumerate}[leftmargin=1.4em,itemsep=2pt]
\item A reproducible description of the fine-tuning architecture
(\S\ref{sec:arch}): static-corpus build, RL-loop online generation,
multi-tier preference data, and the merge/quantise/publish/rollback chain.
\item A chronological benchmark history across versions \texttt{v6}
through \texttt{v14} and the Diamond-SFT and piecewise runs that
followed (\S\ref{sec:runs}), with per-run TLC and SANY counts taken
verbatim from \texttt{outputs/benchmark\_results\_*.csv}.
\item Four engineering lessons distilled from the log
(\S\ref{sec:lessons}): catastrophic forgetting on small corpora, the
binary-reward vacuity trap, the unreliability of quick-eval signal, and
the value of an automated rollback gate.
\item A multi-agent topic-curriculum generator
(\S\ref{sec:diamond-gen-v2}) that produces $200$ Diamond-passing specs
across ten algorithm families at a $200/200$ yield, together with the
holdout-safe SFT corpus and incremental fine-tune that follow from it.
\end{enumerate}

% ---------------------------------------------------------------------------
\section{The Fine-Tuning Architecture}
\label{sec:arch}

The training pipeline is end-to-end automated. The intent is that an
operator can leave the RL loop running unattended for days, and any
checkpoint that survives the eval gate is merged, quantised, packaged as a
GGUF, loaded into the local Ollama daemon, and (optionally) published to
Hugging Face. Anything that fails the eval gate is rolled back. The
sub-systems are described below; readers who want the cheat-sheet view
can skip to Figure~\ref{fig:pipeline}.

\paragraph{Static corpus.} Raw TLA\textsuperscript{+} sources are scraped
into \texttt{data/raw/}, validated with SANY through the
\texttt{pipeline.py}/\texttt{annotate.py} chain, and written to
\texttt{data/validated/combined.jsonl}. A separate dataset builder
(\texttt{src.training.dataset\_builder}) produces task variants for
specification generation, completion, and invariant generation. The
recommended production flags are \texttt{-{}-sany-only}
(drop non-parsing rows), \texttt{-{}-include-augmented} (append the
RL-harvested gold), \texttt{-{}-include-description-sft} (holdout-safe
English-to-spec rows), and \texttt{-{}-bugfix-oversample 2}
(2$\times$ the rows that pair a TLC error with its repair). One subtle
trap: the public corpus is full of \texttt{MC*} wrapper modules whose
only content is \texttt{EXTENDS Real{\_}protocol, TLC}. Training on these
pairs a long English description with a stub specification, which
poisons the spec-generation task; the dataset builder drops them by
default through \texttt{src/training/module\_family.py}.

\paragraph{Online RL loop.} The continuous training loop
(\texttt{scripts/rl\_loop.py}) drives generation, validation, tiering,
augmentation, and retrain decisions. Each cycle has six phases:
(1)~prompt sampling from the benchmark, the synthetic \texttt{RL*}
catalogue, and (optionally) condensed natural-language descriptions for
mapped problems; (2)~spec generation through Ollama against the current
\texttt{chattla:20b} tag, followed by SANY parsing and TLC model
checking; (3)~tiering into gold, silver, bronze, and \emph{bugfix}
(silver $+$ TLC error string $\to$ gold target) classes, with
deduplication by spec hash; (4)~retrain when accumulated new rows cross
a threshold; (5)~quick benchmark every cycle and full benchmark on a
configurable nightly cadence; (6)~difficulty cap from \texttt{rl\_history.jsonl}
so the loop does not over-sample hard prompts while TLC is still weak.
The cycle has no inter-cycle sleep by default.

\paragraph{Retrain chain.} A retrain is a four-step subprocess chain:
\texttt{dataset\_builder} $\to$ \texttt{train.py} (LoRA SFT in harmony
chat format with VRAM-aware \texttt{max\_length}) $\to$
\texttt{merge\_lora.py} (dual-GPU with CPU fallback) $\to$
\texttt{convert\_to\_gguf.py} (\texttt{llama.cpp} convert at \texttt{Q8\_0},
producing \texttt{outputs/gguf/chattla-20b-Q8\_0.gguf} and registering it
locally as \texttt{chattla:20b}). When DPO is enabled, \texttt{train.py
-{}-dpo-after} runs \texttt{train\_dpo.run\_after\_sft} on gold rows in
\texttt{dpo\_pairs.jsonl}, with $\beta = 0.1$ and learning rate
$5\!\times\!10^{-6}$.

\paragraph{Augmented archive and best-per-prompt merging.} The
augmented archive (\texttt{data/processed/augmented.jsonl}) is
append-only across RL cycles. Nothing is deleted as the benchmark moves;
each rebuild instead collapses augmented rows to one row per
\texttt{\_prompt\_id}, keeping the highest-tier survivor (gold and
bugfix dominate silver, which dominates bronze). New gold replaces old
silver for the same prompt. This is the smallest mechanism we found for
keeping the dataset monotonically improving across cycles without
manually pruning it.

\paragraph{Eval gate and rollback.} Every retrain is followed by a small
benchmark (10 problems by default; previously 5). The gate requires at
least three SANY passes and at least one TLC pass; failures trigger an
automatic rollback to the previous GGUF and re-tag the Ollama model.
The gate exists because of an incident in early April 2026 when a
catastrophic SFT run silently shipped a degraded model to Hugging Face;
the gate now sits between training and \texttt{publish\_hf}.

\paragraph{Hugging Face publish gate.}
\texttt{publish\_hf -{}-require-fresh-full-benchmark-hours N} aborts the
publish if no \texttt{outputs/benchmark\_results\_*\_full\_*.csv} is
younger than $N$ hours, and the RL loop reads
\texttt{CHATTLA\_PUBLISH\_REQUIRE\_BENCHMARK\_HOURS} from the
environment. This is the second half of the gate: even a model that
survived the local eval cannot be published unless a recent full
benchmark exists.

\begin{figure}[t]
\centering
\fbox{\parbox{0.95\linewidth}{\small
\textbf{Static:} \texttt{data/raw} $\to$ \texttt{validate} $\to$
\texttt{combined.jsonl} $\to$ \texttt{dataset\_builder} $\to$
\texttt{train.jsonl}, \texttt{eval.jsonl}.\\[2pt]
\textbf{RL cycle:} sample prompts $\to$ Ollama generate $\to$
SANY $+$ TLC $\to$ tier (gold/silver/bronze/bugfix) $\to$ append to
\texttt{augmented.jsonl}, \texttt{dpo\_pairs.jsonl} $\to$
threshold? $\to$ rebuild $+$ retrain.\\[2pt]
\textbf{Retrain:} \texttt{dataset\_builder} $\to$ \texttt{train.py}
(LoRA SFT $\pm$ DPO) $\to$ \texttt{merge\_lora} $\to$
\texttt{convert\_to\_gguf} $\to$ Ollama tag $\to$ \emph{eval gate} $\to$
\texttt{publish\_hf} (or rollback).\\[2pt]
\textbf{Quality:} full-benchmark TLC \% (primary) $\succ$ full-benchmark
SANY \% $\succ$ quick-eval trend $\succ$ structural heuristics
(diagnostic only).
}}
\caption{The end-to-end fine-tuning and deployment pipeline. The
critical property is that no manual step sits between a raw scrape and
a published Hugging Face artefact except for the eval gate.}
\label{fig:pipeline}
\end{figure}

% ---------------------------------------------------------------------------
\section{A Year of Runs}
\label{sec:runs}

\paragraph{The Diamond gate, in two paragraphs.} Several rows below
turn on a semantic-validity criterion we call the \emph{Diamond gate},
defined and audited at length in the methods
paper~\cite{spencer2026chattla} and summarised here so this paper is
self-contained. A specification is Diamond-valid iff it satisfies four
clauses: (1)~SANY accepts the module; (2)~under a small fixed instance,
TLC reaches more than one distinct state; (3)~the module declares at
least one invariant that is neither a syntactic tautology nor the
constant predicate \texttt{TRUE}; and (4)~the spec is mutation-sensitive,
in the sense that a two-phase mutation test (Phase~A negates a guard,
Phase~B negates the leading conjunct of the invariant) produces a TLC
counterexample relative to the unmutated baseline in at least one
phase.

The point of clauses (2)--(4) is to rule out a class of specifications
that pass SANY and TLC by being behaviourally empty: one-state state
spaces, \texttt{UNCHANGED}-dominated next-state relations, decorative
invariants disjoined with \texttt{TRUE}. The methods paper applies the
gate to our internal corpus of $484$ specs previously labelled
gold/silver under SANY$+$TLC and reports that $0$ of them pass; the
present paper takes that audit as a given and uses it to explain why
the preference-optimisation runs in Table~\ref{tab:run-history}
behaved the way they did.

Table~\ref{tab:run-history} records the headline numbers for every
fine-tuning version we ran against the 20-problem held-out benchmark
between roughly January and April 2026. The column conventions are
identical to Table~1 of the methods paper~\cite{spencer2026chattla}; the
underlying CSVs are in
\texttt{outputs/benchmark\_results\_*.csv} and
\texttt{outputs/benchmark\_results/}, and the counts in this table are
taken verbatim from those files.

\begin{table}[t]
\centering
\small
\begin{tabular}{lllcc l}
\toprule
Version & Date & Recipe & SANY & TLC & Outcome \\
\midrule
v6   & 2026-01 & broad SFT, 1 epoch, harmony format         & 5/20 & 1/20 & baseline \\
v7   & 2026-02 & v6 + augmented gold (RL cycles 80--100)     & 6/20 & 2/20 & small lift \\
v8   & 2026-02 & v7 + bugfix oversample 2$\times$            & 7/20 & 2/20 & SANY rises, TLC flat \\
v9   & 2026-03 & best-of-$5$ self-correct at inference       & 7/20 & 2/20 & no train change \\
v10  & 2026-03 & v9 + KTO on 5.7:1 desirable:undesirable    & 4/20 & 1/20 & \textbf{regression} \\
v11  & 2026-03 & v10 reverted, SFT only                     & 6/20 & 2/20 & recovered \\
v12  & 2026-03 & DPO sweep ($\beta\!\in\!\{0.05,0.2\}$)     & ---  & ---  & aborted, see text \\
v13  & 2026-04-03 & v11 + DPO (17 gold pairs, $\beta=0.1$)  & \textbf{9/20} & \textbf{5/20} & best end-to-end \\
v14  & 2026-04-04 & retrain from base, 234 ex, 15 epochs   & 0/20 & 0/20 & \textbf{catastrophic} \\
v14$'$ & 2026-04-04 & v14 SFT-only, 1 epoch                 & 1/20 & 0/20 & still broken \\
v13$_R$ & 2026-04-04 & v13 restored from GGUF backup        & 9/20 & 5/20 & rollback \\
\midrule
diamond-SFT & 2026-04-06 & 720 ex, 73 Diamond, 3 ep, lr 5e-5 & 8/20 & 4/20 & semantic recovery \\
piecewise   & 2026-04-06 & v13$_R$ + 5-step decomposition   & \textbf{12/20} & \textbf{6/20} & inference-time \\
\midrule
\multicolumn{6}{l}{\emph{Diamond-Gen v2 (\S\ref{sec:diamond-gen-v2}), 30-spec held-out diamond eval}} \\
deployed pre & 2026-04-07 & v13$_R$, no RAG, 30 holdout      & ---  & 4/30 (3 dia) & baseline \\
gen-v2 SFT  & 2026-04-07 & 1053 ex (713+170$\times$2), 3 ep, lr 5e-5 & --- & \emph{in flight} & ETA $\sim\!7.5$h \\
\bottomrule
\end{tabular}
\caption{Held-out 20-problem benchmark across model versions. Counts
read directly from the CSVs in \texttt{outputs/}: \texttt{v13} from
\texttt{benchmark\_results\_v13\_full\_20260403.csv} ($9$ SANY, $5$
TLC), \texttt{v14} from \texttt{benchmark\_results\_v14\_full\_20260404.csv}
($0$ SANY, $0$ TLC), \texttt{v14}$'$ from
\texttt{benchmark\_results\_sft\_only\_20260404.csv} ($1$ SANY, $0$
TLC), \texttt{diamond-SFT} from
\texttt{benchmark\_results/benchmark\_results\_diamond\_sft\_20260406\_002008.csv}
($8$ SANY, $4$ TLC), and the piecewise row from
\texttt{benchmark\_results/piecewise\_first\_run\_20260406.csv} ($12$
SANY, $6$ TLC) which uses the same underlying \texttt{v13}-restored
checkpoint as its single-shot baseline. The bottom block reports the
two Diamond-Gen v2 measurements (\S\ref{sec:diamond-gen-v2}) which are
on a separate $30$-spec held-out diamond benchmark and are not directly
comparable to the $20$-problem TLC counts above; the
\texttt{deployed pre} row is the pre-train baseline measured before
launching the v2 SFT and is read directly from
\texttt{outputs/eval/holdout\_pre.json} ($4$ gold of which $3$ pass the
Diamond gate, $11$ silver, $15$ bronze).}
\label{tab:run-history}
\end{table}

\paragraph{The shape of the curve.} Reading down the table, three
features stand out. (1)~The reward of incremental SFT $+$ DPO peaks at
\texttt{v13}, with $5/20$ TLC. (2)~The single largest delta in either
direction is the \texttt{v14} catastrophe: a from-scratch SFT pass on
$234$ examples for $15$ epochs at learning rate $3\!\times\!10^{-4}$
zeroed both SANY and TLC, which is a striking failure given that the
same examples had been part of \texttt{v13}'s training set without
incident. The difference was the schedule, not the data. (3)~The
piecewise inference procedure on the restored \texttt{v13} checkpoint
adds $+3$ SANY and $+1$ TLC \emph{without} touching the weights, which
is a stronger per-FLOP gain than any of the training interventions
above.

\paragraph{What happened to v12.} The version label \texttt{v12} is
intentionally absent from the analysis above and corresponds to a DPO
hyperparameter sweep ($\beta \in \{0.05, 0.2\}$, two LR settings) that
we ran in late March against the same $17$-pair gold set later used by
\texttt{v13}. The sweep aborted before producing a publishable
checkpoint: the $\beta=0.05$ run diverged early in training, the
$\beta=0.2$ run produced reward margins indistinguishable from the
\texttt{v11} baseline on the canary set, and we cancelled the sweep
in favour of the single $\beta=0.1$ configuration that became
\texttt{v13}. There is no full benchmark CSV for \texttt{v12} because
no checkpoint from the sweep cleared the eval gate. We mention this
explicitly so the version numbering is not mysterious; nothing in
\texttt{v12} contradicts the rest of the table.

\paragraph{The KTO regression.} The KTO run (\texttt{v10}) is worth a
short autopsy because it taught us the most. We had accumulated
several hundred gold/silver pairs from RL cycles and assembled a KTO
training set with a desirable-to-undesirable ratio of $5.7\!:\!1$.
After training, the held-out TLC count dropped from $2/20$ to $1/20$
and the canary regression set fell to $0/5$. The verifier metrics that
drove the preference labels were entirely binary (\texttt{tlc\_pass}),
and the gold tier turned out to contain a large fraction of vacuously
correct specifications: parser-clean, model-checker-clean, and
behaviourally empty. The Diamond audit~\cite{spencer2026chattla} is the
follow-on diagnosis: $0$ of $484$ specs in the gold/silver pool passed
mutation testing. KTO faithfully maximised the binary signal it was
asked to maximise; the signal was wrong.

\paragraph{The \texttt{v14} catastrophe.} \texttt{v14} was a regression
test of the dataset builder under the new DPO-first plumbing. The
intent was a clean retrain over the same data \texttt{v13} had used.
The actual run trained from the unmodified \texttt{gpt-oss:20b} base
weights for $15$ epochs over $234$ examples at learning rate
$3\!\times\!10^{-4}$. SANY pass dropped from $9/20$ to $0/20$. A
follow-up ``SFT-only, 1 epoch'' experiment (\texttt{v14}$'$) recovered
exactly one SANY pass and zero TLC passes; the model had forgotten how
to write TLA\textsuperscript{+} entirely. Recovery required restoring
the previous \texttt{v13} GGUF from backup and re-tagging Ollama, which
the eval gate (\S\ref{sec:arch}) had been built to do.

\paragraph{The Diamond-SFT recovery.} The Diamond-SFT run on
2026-04-06 used $720$ examples assembled as $73$ Diamond-passing specs
oversampled $3\!\times$, $205$ gold RL specs, $231$ gold benchmark
specs, and $86$ description-SFT rows, trained for $3$ epochs at
learning rate $5\!\times\!10^{-5}$ with the merged \texttt{v13}
checkpoint as the base (incremental, not from scratch). Training loss
fell from $1.72$ to $0.49$. The benchmark recovered to $8/20$ SANY and
$4/20$ TLC---below \texttt{v13}'s peak in TLC but with much higher
structural quality ($0.92$ versus $0.77$ average), and with the
remaining failures cleanly attributable to small SANY syntax mistakes
(\texttt{EXISTS} instead of \texttt{\textbackslash E}, \texttt{SUM} as
an undeclared operator, \texttt{EXCEPT} comma syntax, \texttt{CHOOSE}
parenthesisation). This is a more useful failure mode than the
\texttt{v13} one: structurally-correct specs with fixable parser
errors are exactly what an automated SANY repair pass would close.

\paragraph{RL cycle volume.} For context on the scale, the RL loop ran
through approximately cycle~\texttt{c266} during the period covered by
this paper, with full benchmark CSVs captured at every nightly cadence
(see \texttt{outputs/benchmark\_results/benchmark\_results\_rl\_c*\_full\_*.csv}).
The cycle-to-cycle variance on the quick eval is several problems,
which is why the table above only reports full benchmarks; the quick
eval is best read as a trend signal, not as a publication-quality
metric.

% ---------------------------------------------------------------------------
\section{Four Lessons from the Log}
\label{sec:lessons}

\paragraph{Lesson 1: Do not train broadly from scratch on small data.}
The \texttt{v14} catastrophe is the cleanest example of catastrophic
forgetting we have ever seen on this codebase. The $234$ examples
involved had all appeared in \texttt{v13} without incident; the
difference was that \texttt{v13} added them on top of an already
TLA\textsuperscript{+}-fluent merged checkpoint, while \texttt{v14}
re-derived the adapter from the unmodified base. Fifteen epochs at
$3\!\times\!10^{-4}$ over a few hundred examples is enough to overwrite
the prior. The operational rule we adopted afterwards is that the SFT
recipe always builds on the previous merged checkpoint, learning rate
is capped at $1\!\times\!10^{-5}$ for incremental runs, and epoch count
is hard-capped at $2$ in the loop driver.

\paragraph{Lesson 2: Binary verifier rewards are an attractor towards
vacuous specs.} KTO and DPO are both preference-learning losses; both
are well-understood; both behaved correctly given the labels we gave
them. The labels were the bug. \texttt{tlc\_pass} is satisfied by
\emph{any} specification whose state space contains no invariant
violations, and the easiest way to satisfy that is to write a
specification with a one-state state space and a tautological
invariant. The training corpus we built up over months of RL was
densely populated with exactly that class of spec
($0/484$ Diamond-passing under the audit). Any preference-optimisation
loop pointed at this corpus is not learning ``be correct''; it is
learning ``be small and quiet.'' The fix is upstream: define
correctness as a four-clause Diamond gate that includes mutation
sensitivity, harvest the corpus by rejection sampling against that
gate, and only \emph{then} run the preference loss. The full
description of the gate is in our methods
paper~\cite{spencer2026chattla}.

\paragraph{Lesson 3: The quick eval is a trend, not a metric.} The
default quick benchmark is $12$ problems. Across cycles
\texttt{c100}--\texttt{c260}, the quick eval moved by one or two
problems essentially every cycle, in either direction, with no
correlation to actual model improvements. The full benchmark
($20$ problems) is more stable but still has $\pm 1$--$2$ problem
wobble across reruns. The operational rule we adopted is: never
report a result based on the quick eval, never make a deploy/rollback
decision on a single full eval, and require a fresh full eval within
$N$ hours before allowing the publish step to proceed. Half the
publish-related code in the loop exists to enforce this.

\paragraph{Lesson 4: Build the rollback before you need it.} The
\texttt{v14} regression caused real damage in the sense that a
degraded model was, briefly, the published \texttt{chattla:20b} tag.
The eval gate, the ``require fresh full benchmark'' flag on
\texttt{publish\_hf}, and the GGUF backup-and-restore step were all
written in response to that incident. None of them are interesting
ML; all of them are interesting MLOps. The asymmetry between the
cost of building rollback infrastructure (one afternoon) and the cost
of not having it (a full day of restoration plus a published
regression) is severe enough that we now treat rollback as a
prerequisite, not a nice-to-have.

\begin{table}[t]
\centering
\small
\begin{tabular}{lll}
\toprule
Symptom & Recipe that caused it & Recipe that fixed it \\
\midrule
TLC $\to 0$, SANY $\to 0$ & broad SFT from base, $15$ epochs, lr $3e\!-\!4$ &
  rollback to merged \texttt{v13} GGUF \\
gold tier full of vacuous specs & binary \texttt{tlc\_pass} reward &
  Diamond gate $+$ rejection sampling \\
quick-eval lurches every cycle & $12$-problem subset & full $20$ benchmark \\
silent regression in production & no eval gate &
  3-SANY + 1-TLC eval gate, $N$-hour fresh-bench gate \\
\bottomrule
\end{tabular}
\caption{Failure modes encountered in the fine-tuning log and the
operational fix that addressed each one. None of the right-hand
column changes touch the model architecture.}
\label{tab:lessons}
\end{table}

% ---------------------------------------------------------------------------
\section{Diamond-Filtered SFT}
\label{sec:diamond-sft}

The current recipe, in the wake of all of the above, is small and
boring. We restore the merged \texttt{v13} checkpoint, audit every
candidate training example through the Diamond gate
(SANY pass, more than one reachable state, declared non-vacuous
invariant, mutation-sensitive), oversample the survivors $3\!\times$,
mix in $86$ description-SFT rows for English-to-spec coverage, and
train incrementally for $3$ epochs at $5\!\times\!10^{-5}$ on top of
the existing weights. The rejection sampler lives in
\texttt{scripts/diamond\_sft\_gen.py}, which uses a teacher model
(\texttt{gpt-oss:120b} via Ollama) to propose candidates and then
discards anything that fails the gate. Of the $484$ specs in the
historical augmented archive, the gate keeps $73$. Of the new
candidates the teacher generates, the keep rate is similar---the
binding constraint is not access to a teacher, it is the mutation
clause.

The result of running this recipe is the \texttt{diamond-SFT} row in
Table~\ref{tab:run-history}: $8/20$ SANY, $4/20$ TLC, structural score
$0.92$. That is one TLC pass below \texttt{v13}, but with a much
healthier failure profile---the residual SANY failures are local
syntax slips, not structural collapse---and with a corpus that the
KTO/DPO branches can be safely re-enabled against. Combined with the
piecewise inference procedure of the methods
paper~\cite{spencer2026chattla}, the same \texttt{v13}-derived
checkpoint reaches $12/20$ SANY and $6/20$ TLC. Our current
operating position is therefore: train incrementally on
Diamond-filtered data, generate at inference time through piecewise
decomposition, and reserve preference optimisation for the narrow
follow-on case where the Diamond corpus is dense enough to support
it.

% ---------------------------------------------------------------------------
\section{Discussion}

\paragraph{Why broad SFT keeps failing.} The TLA\textsuperscript{+}
public corpus is tiny by modern standards (a few hundred specs), and
the canonical examples are short, idiomatic, and structurally similar.
A $20$B model already has a passable prior over this distribution from
pre-training, so the marginal value of broad SFT against more of the
same is small. What broad SFT \emph{does} do reliably is overwrite
adjacent parts of the prior---braces, identifiers, English idioms---in
ways that show up as catastrophic forgetting on benchmark problems
where the model would otherwise have stitched together knowledge from
non-TLA\textsuperscript{+} parts of its training. The one regime in
which we got SFT to help was Diamond-filtered, incremental, low-LR,
short-epoch.

\paragraph{Why preference optimisation against a binary verifier
fails.} A preference loss only has access to the labels you give it. A
TLA\textsuperscript{+} verifier emits a binary acceptance signal that
is satisfied by an entire equivalence class of specifications, only
some of which a human reviewer would call correct. Optimising hard
against the binary signal moves the model towards the easy-to-reach
corners of that class (degenerate state spaces, decorative
invariants), which is exactly the regime the methods paper documents
through the Diamond gate. The fix is to make the reward semantically
sharper before training, not to make the optimiser stronger.

\paragraph{Why inference-time decomposition is the cheapest win.}
Decomposition costs roughly $5\!\times$ the per-spec inference budget
relative to single-shot, but it requires no training data, no
labelling, no rollback infrastructure, and no risk to the published
checkpoint. In a regime where each round of fine-tuning has roughly
even odds of helping or hurting, ``do not train, generate more
carefully'' is a reasonable default. We do not believe inference-time
methods replace fine-tuning---the Diamond-SFT recipe is what gives the
piecewise procedure something to decompose---but the order matters.
Train incrementally on small, semantically audited corpora; generate
through verifier-guided decomposition; treat broad fine-tuning as a
last resort.

\paragraph{Diamond-Gen v2 in one paragraph.} The current
operating position above was written before the
\S\ref{sec:diamond-gen-v2} round, in which a multi-agent
topic-curriculum generator produced $200$ Diamond-passing specs at a
$200/200$ yield and the resulting incremental SFT is at the time of
writing in flight. The Diamond-Gen v2 result does not change the
Diamond-Filtered SFT recipe described above; it changes the
\emph{source} of the Diamond corpus from teacher rejection sampling at
$\sim\!15\%$ yield to a directed generator at $100\%$, and removes
``access to a gate-passing corpus'' from the list of binding
constraints.

\paragraph{Related work in two paragraphs.} Our negative results
echo a long line of work on catastrophic forgetting in fine-tuning
\cite{kirkpatrick2017ewc,lopez2017gem}: even moderate-LR SFT on a
narrow distribution will overwrite adjacent capabilities of a base
model, and the standard mitigations---elastic weight consolidation,
gradient episodic memory, replay buffers---all encode the same
intuition we eventually adopted by hand, namely that incremental
training on top of a previous merged checkpoint with a low LR and
small epoch budget is safer than re-deriving the adapter from base.
We do not implement EWC or replay explicitly; the augmented archive's
best-per-prompt merge (\S\ref{sec:arch}) plays a similar role at the
data level by ensuring that previously-seen prompts remain present
across rebuilds.

For fine-tuning on formal-language targets specifically, the closest
prior work we are aware of is on Dafny~\cite{misu2024towards} and
Isabelle/Coq~\cite{first2023baldur}; both report that broad SFT on
small formal corpora is fragile in ways that mirror our experience,
and both end up relying on verifier feedback as a selection mechanism
rather than as a gradient. Our companion methods
paper~\cite{spencer2026chattla} situates the
TLA\textsuperscript{+}-specific picture against the systematic
$30$-model evaluation in~\cite{spencer2026canllms}, which reports a
$8.6\%$ semantic-correctness ceiling under SANY$+$TLC across $2{,}730$
runs and is the broadest baseline measurement available for this
language.

\paragraph{Statistical caveats on $N=20$.} The held-out benchmark has
$20$ problems. Treating each problem as an independent Bernoulli
trial, a Wilson $95\%$ confidence interval on a point estimate of
$5/20$ is approximately $[0.11, 0.47]$, which is wide enough that the
$5/20$ \texttt{v13} TLC count and the $4/20$ Diamond-SFT count are not
statistically distinguishable; the $6/20$ piecewise count and the
$1/20$ baseline used to size the piecewise gain in the methods paper
\emph{are} distinguishable at $\alpha=0.05$ under the same
approximation, but only barely. The right reading of
Table~\ref{tab:run-history} is therefore: the catastrophic deltas
(\texttt{v14} $0/0$, KTO regression to $1/20$ TLC) are large enough
that no plausible CI overlaps the \texttt{v13} peak, and we treat
those as load-bearing signal; the small deltas between
adjacent rows (\texttt{v8} versus \texttt{v9}, \texttt{v13} versus
diamond-SFT) sit inside the wobble and should be read as
\emph{directional} only. Conclusions in \S\ref{sec:lessons} that
depend on the catastrophic deltas (Lessons 1, 2, 4) are robust to the
small $N$; the Lesson 3 claim about quick-eval volatility is observed
across $\sim\!200$ cycles and does not depend on benchmark size.

\paragraph{Limitations.}
\emph{One model.} Every result above uses one underlying checkpoint,
\texttt{gpt-oss:20b}. We do not yet know whether the catastrophic
forgetting threshold (LR, epoch count, sample count) generalises to
larger or coding-specialised models, or whether the binary-reward
vacuity trap is just as bad with a stronger prior. The conjecture is
that both effects soften with scale but do not vanish.
\emph{One benchmark.} The 20-problem held-out suite is small and
self-authored. The methods paper~\cite{spencer2026chattla} discusses
this at length and we do not duplicate that discussion here.
\emph{No multi-seed.} Every cell in Table~\ref{tab:run-history} is a
single run. Informal repeats suggest a $\pm 1$--$2$ problem wobble.
The deltas we treat as load-bearing (\texttt{v14} catastrophe,
piecewise gain, KTO regression) are larger than the wobble; the
deltas we treat as suggestive (\texttt{v8} versus \texttt{v9}) are
not.

% ---------------------------------------------------------------------------
\section{Diamond-Gen v2: Topic-Curriculum Multi-Agent Generation}
\label{sec:diamond-gen-v2}

The Diamond-Filtered SFT recipe of \S\ref{sec:diamond-sft} is
bottlenecked by the keep rate of its rejection sampler. Out of $484$
specs in the augmented archive, $73$ pass the Diamond gate; the keep
rate on freshly proposed teacher candidates is similar. Improving the
recipe therefore turns on either (a) raising the rejection-sampler
keep rate or (b) replacing the rejection sampler with a directed
generator. The round documented in this section takes the second path.

\paragraph{Topic curriculum.} We hand-curate a $200$-topic seed list
($10$ algorithm families $\times$ $20$ topics each) covering
concurrency primitives, classical mutual-exclusion algorithms,
consensus and election protocols, point-to-point and broadcast
communication, bounded data structures, transactional and replicated
storage, scheduling and resource allocation, memory hierarchies and
cache coherence, workflow and state-machine FSMs, and classical
puzzles. Each topic specifies a target module name, a one-paragraph
problem description, and (where applicable) the canonical algorithm
the spec must implement faithfully (e.g.\ \emph{Peterson, two-process,
flag and turn}; \emph{Lamport bakery, choosing\,/\,number arrays};
\emph{single-decree Paxos, prepare\,/\,promise\,/\,accept ballots}). The
seed list lives in \texttt{data/diamond\_gen\_topics.json}.

\paragraph{Per-batch agents.} We dispatch the $10$ batches in parallel
to $10$ instances of a strong generation model (Claude Opus 4.6, $1$M
context, called as Claude Code subagents). Each agent receives only
its own $20$-topic batch, the path to a thin validator CLI
(\texttt{scripts/validate\_diamond\_cli.py}) that wraps
\texttt{validate\_diamond} from
\texttt{scripts/diamond\_sft\_gen.py}, and a strict prompt encoding
(i)~the four Diamond clauses, (ii)~a required spec skeleton (file
basename = module, single \texttt{TypeOK} that conjoins the strong
safety property), (iii)~a hard $\le\!200$ reachable-state ceiling so
TLC clears the $60\,\text{s}$ timeout, and (iv)~a per-topic budget of
up to six retry-on-failure attempts. The agent runs an inner loop:
write spec $\rightarrow$ shell out to the validator $\rightarrow$ read
the JSON verdict $\rightarrow$ on \texttt{bronze} fix the SANY error,
on \texttt{silver} shrink the constants or add a fairness condition,
on \texttt{gold-not-diamond} strengthen the safety conjunct or expand
the action set. The agent emits a $20$-record JSONL when its batch
completes.

\paragraph{Independent re-validation.} Once all ten batches return,
\texttt{scripts/aggregate\_diamond\_gen.py} re-runs the Diamond
validator on every spec from a clean process pool ($8$ workers) and
writes the authoritative aggregate
\texttt{outputs/diamond\_gen/\\diamond\_generated.jsonl}. The clean
re-validation deliberately does not trust the per-agent self-reports.
On the $200$-spec round documented here, the re-validation completed
in $68\,\text{s}$ and the result was $200/200$ at gold tier with all
four Diamond clauses satisfied. State-space sizes ranged from
$4$ (TricolorGc) to $181{,}440$ (FifteenPuzzle, $3{\times}3$ blank);
median $\sim\!42$. No batch had any failures.

\paragraph{Held-out eval split.} The Diamond-Filtered SFT round of
\S\ref{sec:diamond-sft} did not carve a held-out eval set, which is
fine for benchmarking against the $20$-problem suite of
\S\ref{sec:runs} but does not let us measure whether the new
generator is teaching the smaller fine-tuned model anything beyond the
specs it has seen. The v2 round therefore deterministically holds out
the first three specs (sorted by module name) from each of the ten
batches, giving a $30$-spec diamond eval at
\texttt{data/processed/diamond\_eval\_holdout.jsonl}. The carve script
(\texttt{scripts/carve\_diamond\_holdout.py}) loads the existing
\texttt{train.jsonl} first and excludes any module that already
appears there from the holdout pool, which caught \texttt{Barrier} as
a leak (\texttt{Barrier} was already in the v13 training set, so its
pre-train baseline would have been biased) and rotated it for
\texttt{CountDownLatch}. The remaining $170$ generated specs go into
the SFT training pool.

\paragraph{Hand-written reasoning rationales.} The Diamond-Filtered
SFT records of \S\ref{sec:diamond-sft} have a four-turn chat structure
(developer prompt $\rightarrow$ user request $\rightarrow$ assistant
design rationale $\rightarrow$ assistant spec). The rationale turn is
load-bearing: it teaches the student model to plan the variable set
and the safety argument before emitting the module. Generating
$170$ rationales by template extraction does not work---a first attempt
that string-formatted rationales from spec metadata produced $170$
records with the same skeleton, which would teach the student to
produce that skeleton. We therefore dispatch the rationale generation
as four parallel agents with strict no-template prompts and $43$
or $42$ modules each, with a per-batch sanity check that no two
rationales share their first $60$ characters and that every rationale
falls in a $600$--$1500$-character window. The four batches return
$170$ unique-prefix rationales (mean length $1039$, range $813$--$1287$)
that read as per-spec design memos rather than as filled-in templates.

\paragraph{Mixed SFT corpus.} The final training corpus
(\texttt{data/processed/diamond\_sft\_v3.jsonl}, built by
\texttt{scripts/build\_diamond\_sft\_v3.py}) consists of the existing
$713$-record \texttt{train.jsonl} that the v13 checkpoint was trained
on, plus the $170$ new diamond-gen-v2 specs in chat-message format
oversampled $2\!\times$, for $1053$ records total. The build script
hard-asserts that no record contains any of the $30$ holdout module
names anywhere in its assistant turns; running the same check
externally on the final file confirms zero leaks. Oversampling $2\!\times$
is a deliberate compromise: pure incremental training on only the
$170$ new specs risks catastrophic forgetting on the categories the
$713$-record base covers, while training on the union without
oversampling dilutes the gradient contribution of the new
algorithm-family diversity below $20\%$ of the corpus.

\paragraph{Incremental SFT recipe.} The fine-tune is run on top of the
deployed merged checkpoint (\texttt{outputs/merged\_model}, which
\emph{is} v13 + checkpoint-131 LoRA already merged), with a fresh
LoRA adapter, $3$ epochs, learning rate $5\!\times\!10^{-5}$, max
sequence length $4096$, effective batch size $8$ (per-device $2$ over
two GPUs $\times$ $4$ gradient accumulation), bf16, on the same
$2{\times}$RTX 8000 hardware as the rest of the paper. The
\texttt{--resume outputs/checkpoints/checkpoint-131} path that one
would normally use to keep the same LoRA trajectory is broken on
PyTorch 2.6 because the saved \texttt{rng\_state.pth} contains numpy
globals that the new \texttt{weights\_only=True} default in
\texttt{torch.load} refuses to unpickle; the
\texttt{--base-model outputs/merged\_model} path is functionally
equivalent (same starting weights) and avoids the resume-state
machinery entirely. The total schedule is $396$ optimisation steps,
estimated wall clock $\sim\!7.5\,\text{h}$ on the trainer's
bf16-dequantised gpt-oss MoE.

\paragraph{Pre-train baseline on the held-out diamond eval.} Before
launching the SFT, we measure the deployed v13$_R$ checkpoint on the
$30$-spec held-out diamond eval with RAG disabled. The script
(\texttt{scripts/eval\_diamond\_holdout.py}) prompts the deployed
model with each topic description, runs the generated module through
the same Diamond validator used at corpus build time, and tallies
tier counts. The baseline result is $4/30$ at gold tier of which
$3/30$ pass the full Diamond gate---i.e.\ $10\%$ Diamond pass rate on
problems the deployed model has not seen---against $11/30$ silver
(SANY-clean but TLC could not run) and $15/30$ bronze (SANY-fail).
This is the bar the v2 fine-tune has to beat to be deployed as a
\texttt{chattla:20b-v2} replacement for the live tag.

\paragraph{Decision rule.} If the post-train v2 model on the same
$30$-spec eval beats $3/30$ Diamond by a clear margin (we use
``$\ge\!8/30$, with the $11$ baseline-silver problems breaking
predominantly toward gold rather than toward bronze'' as the
operational definition), the new adapter is merged through the
existing
\texttt{deploy\_diamond\_model.sh} pipeline and re-tagged
\texttt{chattla:20b-v2}. If it does not, the corpus is not the
binding constraint and the next lever is the recipe (LR schedule,
oversample weight, max sequence length, fairness conjuncts in the
training rationale) rather than more data. The post-train comparison
will be filled into Table~\ref{tab:run-history} when the run completes;
the in-flight \texttt{gen-v2 SFT} row records the schedule, the corpus
size, and the pre-train baseline so a reader of this revision can
reproduce the experiment without waiting for the post-train number.

\paragraph{What the new method changes about the binding constraint.}
The Diamond-Filtered SFT recipe of \S\ref{sec:diamond-sft} is
constrained by access to a teacher model strong enough to propose
gate-passing candidates at any reasonable rate, and---once such a
teacher exists---by the proportion of its outputs that the gate
accepts. Both constraints translate into a fixed pool size, around
$73$ specs in the round documented in \S\ref{sec:diamond-sft}, plus
whatever else the teacher generates between rebuilds. The Diamond-Gen
v2 method removes both constraints: the topic curriculum tells the
generator \emph{what} to write rather than waiting for the gate to
filter what comes back, and the per-topic retry budget lets the
generator iterate on its own validator output until each spec passes
or a fixed budget is exhausted. The pool size becomes a function of
how many topics the curators can write down rather than how many
candidates a teacher can generate, and the binding constraint shifts
from \emph{teacher access} to \emph{topic coverage and per-agent retry
budget}. Empirically, on the $200$-topic round documented here, the
gate-pass yield was $200/200$ on the first run, the wall clock to
generate and re-validate the entire corpus was approximately one hour
across ten parallel agents, and no topic exhausted its six-retry
budget.

\paragraph{Why retry inside the generator instead of outside it.}
There is a tempting middle path in which a single generator emits
many candidates per topic and the validator filters them downstream,
which is closer to the rejection-sampler shape of
\S\ref{sec:diamond-sft}. We do not take that path because the
information that the validator returns is rich---tier, distinct-state
count, mutation outcome, specific SANY error string---and feeding it
\emph{back into the generator} as input to its next attempt is
strictly more informative than treating each candidate as
independent. Empirically, the agents that get a \texttt{tier=silver}
verdict on attempt $1$ shrink their constants on attempt $2$; the
agents that get \texttt{trivial\_invariant=true} strengthen the
safety conjunct on attempt $2$; the agents that get a SANY parse
error correct the specific token on attempt $2$. Most topics resolved
within two attempts; the longest chains we observed used four. A
rejection sampler that discards failed candidates would have to
re-discover the same fixes from scratch on every retry.

\paragraph{Limitations of the v2 method.} Three caveats apply.
\emph{First}, the generator is a single underlying model. Every spec
in the v2 corpus reflects Claude Opus 4.6's idiom (history-flag
encoding for transitive safety properties, \texttt{TypeOK} as the
single named invariant the auto-generated cfg picks up,
\texttt{Done}\,/\,\texttt{Reset} self-stutter actions to dodge TLC's
default deadlock check). The student model trained on this corpus
will inherit those idioms; whether they generalise to TLA+ written in
other styles is not yet measured. \emph{Second}, the topic curriculum
is hand-curated, which is a labour cost the rejection-sampler
methodology does not have. Scaling beyond $200$ specs without
diminishing returns will likely require either a larger curated topic
list (cheap but bounded by curator effort) or a generator that
proposes new topics by perturbing existing ones in
algorithm-faithful ways (an open question). \emph{Third}, the
per-topic budget cap of six attempts is arbitrary; we set it high
enough that no topic in the $200$-spec round hit the cap, but a
harder topic distribution might. We do not yet know the failure mode
of an agent that exhausts its budget in the middle of a batch.

% ---------------------------------------------------------------------------
\section{Conclusion}

A year of fine-tuning a $20$B open model for TLA\textsuperscript{+}
specification generation taught us that the productive interventions
were not the ones we expected. Broad SFT failed almost every time we
tried it, in proportion to how broad it was. Preference optimisation
against a binary verifier reward drifted the model towards
semantically vacuous specifications. The pipeline-engineering work
that mattered most was the eval gate and the rollback path, which
prevented a single bad retrain from becoming a published regression.
The single largest improvement on the held-out benchmark came from a
five-step inference-time decomposition that does not touch the
weights at all. The current recipe---incremental SFT on
Diamond-filtered data, plus piecewise inference---is small, boring,
and reproducible, and that is the point. For low-resource formal
languages with deterministic verifiers, the leverage is in the shape
of the reward and the structure of the inference, not in the volume
of supervised data.

The April 7 revision adds one further refinement
(\S\ref{sec:diamond-gen-v2}): when the Diamond pool itself is the
binding constraint, a topic-curriculum multi-agent generator that
calls the validator from inside its own retry loop produces
gate-passing specs at $200/200$ yield rather than $\sim\!15\%$, shifts
the binding constraint from teacher access to topic coverage, and
gives the incremental fine-tune a $\sim\!4\!\times$ larger
algorithm-family-diverse corpus to train on. Whether the resulting
post-train checkpoint actually clears the held-out diamond eval by a
margin worth deploying is the open question that the
\texttt{gen-v2 SFT} run in Table~\ref{tab:run-history} will answer
when it lands.

% ---------------------------------------------------------------------------
\bibliographystyle{plain}
\begin{thebibliography}{99}

\bibitem{lamport2002specifying}
Leslie Lamport.
\newblock \emph{Specifying Systems: The TLA\textsuperscript{+} Language and
Tools for Hardware and Software Engineers}.
\newblock Addison-Wesley, 2002.

\bibitem{spencer2026chattla}
Eric Spencer.
\newblock ChatTLA: Verifier-guided decomposition for
TLA\textsuperscript{+} specification generation.
\newblock Manuscript, 2026.

\bibitem{spencer2026canllms}
Eric Spencer.
\newblock Can LLMs write correct TLA\textsuperscript{+} specifications?
Evaluating natural-language-to-TLA\textsuperscript{+} generation.
\newblock Manuscript, 2026.

\bibitem{rafailov2023direct}
Rafael Rafailov, Archit Sharma, Eric Mitchell, Stefano Ermon,
Christopher D. Manning, and Chelsea Finn.
\newblock Direct preference optimization: Your language model is
secretly a reward model.
\newblock In \emph{NeurIPS}, 2023.

\bibitem{ethayarajh2024kto}
Kawin Ethayarajh, Winnie Xu, Niklas Muennighoff, Dan Jurafsky, and
Douwe Kiela.
\newblock KTO: Model alignment as prospect theoretic optimization.
\newblock arXiv:2402.01306, 2024.

\bibitem{hu2022lora}
Edward J. Hu, Yelong Shen, Phillip Wallis, Zeyuan Allen-Zhu, Yuanzhi
Li, Shean Wang, Lu Wang, and Weizhu Chen.
\newblock LoRA: Low-rank adaptation of large language models.
\newblock In \emph{ICLR}, 2022.

\bibitem{kirkpatrick2017ewc}
James Kirkpatrick, Razvan Pascanu, Neil Rabinowitz, Joel Veness,
Guillaume Desjardins, Andrei A. Rusu, Kieran Milan, John Quan,
Tiago Ramalho, Agnieszka Grabska-Barwinska, Demis Hassabis, Claudia
Clopath, Dharshan Kumaran, and Raia Hadsell.
\newblock Overcoming catastrophic forgetting in neural networks.
\newblock \emph{Proceedings of the National Academy of Sciences},
114(13):3521--3526, 2017.

\bibitem{lopez2017gem}
David Lopez-Paz and Marc'Aurelio Ranzato.
\newblock Gradient episodic memory for continual learning.
\newblock In \emph{NeurIPS}, 2017.

\bibitem{misu2024towards}
Md Rakib Hossain Misu, Cristina V. Lopes, Iris Ma, and James Noble.
\newblock Towards AI-assisted synthesis of verified Dafny methods.
\newblock \emph{Proceedings of the ACM on Software Engineering},
1(FSE), 2024.

\bibitem{first2023baldur}
Emily First, Markus N. Rabe, Talia Ringer, and Yuriy Brun.
\newblock Baldur: Whole-proof generation and repair with large
language models.
\newblock In \emph{ESEC/FSE}, 2023.

\end{thebibliography}

\end{document}
